\begin{thebibliography}{00}

\bibitem{b_llmedu} L. Yan \emph{et al.}, ``Practical and ethical challenges of large language models in education: A systematic scoping review,'' \emph{British Journal of Educational Technology}, vol. 55, no. 1, pp. 90--112, 2024.

\bibitem{b_gpt3} T. Brown \emph{et al.}, ``Language models are few-shot learners,'' in \emph{Advances in Neural Information Processing Systems}, vol. 33, 2020, pp. 1877--1901.

\bibitem{b_survey_edu} X. Zhai \emph{et al.}, ``Large language models in education: Vision and opportunities,'' \emph{arXiv preprint arXiv:2311.13160}, 2023.

\bibitem{b_lpp} X. Wang \emph{et al.}, ``Educational personalized learning path planning with large language models,'' \emph{arXiv preprint arXiv:2407.11773}, 2024.

\bibitem{b_ramo} J. Zhang \emph{et al.}, ``RAMO: Retrieval-augmented generation for enhancing MOOCs recommendations,'' \emph{arXiv preprint arXiv:2407.04925}, 2024.

\bibitem{b_eduplanner} X. Zhang \emph{et al.}, ``EduPlanner: LLM-based multi-agent systems for customized and intelligent instructional design,'' \emph{IEEE Transactions on Learning Technologies}, 2025, doi: 10.1109/TLT.2025.3555880.

\bibitem{b_masla} T. Guo \emph{et al.}, ``Large language model based multi-agents: A survey of progress and challenges,'' in \emph{Proc. 33rd Int. Joint Conf. Artificial Intelligence (IJCAI)}, 2024.

\bibitem{b_react} S. Yao \emph{et al.}, ``ReAct: Synergizing reasoning and acting in language models,'' in \emph{Proc. Int. Conf. Learning Representations (ICLR)}, 2023.

\bibitem{b_reflexion} N. Shinn, F. Cassano, A. Gopinath, K. Narasimhan, and S. Yao, ``Reflexion: Language agents with verbal reinforcement learning,'' in \emph{Advances in Neural Information Processing Systems}, vol. 36, 2023.

\bibitem{b_selfrefine} A. Madaan \emph{et al.}, ``Self-Refine: Iterative refinement with self-feedback,'' in \emph{Advances in Neural Information Processing Systems}, vol. 36, 2023.

\bibitem{b_langgraph} LangChain Inc., ``LangGraph: Building stateful, multi-actor applications with LLMs,'' 2024. [Online]. Available: https://langchain-ai.github.io/langgraph/

\bibitem{b_judge} L. Zheng \emph{et al.}, ``Judging LLM-as-a-judge with MT-Bench and Chatbot Arena,'' in \emph{Advances in Neural Information Processing Systems}, vol. 36, 2023.

\bibitem{b_pydantic} Pydantic, ``Pydantic: Data validation using Python type hints,'' 2024. [Online]. Available: https://docs.pydantic.dev/

\bibitem{b_halluc} Z. Ji \emph{et al.}, ``Survey of hallucination in natural language generation,'' \emph{ACM Computing Surveys}, vol. 55, no. 12, pp. 1--38, 2023.

\bibitem{b_embed} OpenAI, ``New embedding models and API updates,'' 2024. [Online]. Available: https://openai.com/index/new-embedding-models-and-api-updates/

\bibitem{b_lp_survey} A. H. Nabizadeh, J. P. Leal, H. N. Rafsanjani, and R. R. Yadav, ``Learning path personalization and recommendation methods: A survey of the state-of-the-art,'' \emph{Expert Systems with Applications}, vol. 159, p. 113596, 2020.

\bibitem{b_rag} P. Lewis \emph{et al.}, ``Retrieval-augmented generation for knowledge-intensive NLP tasks,'' in \emph{Advances in Neural Information Processing Systems}, vol. 33, 2020, pp. 9459--9474.

\bibitem{b_agents_matter} S. Kapoor, B. Stroebl, Z. S. Siegel, N. Nadgir, and A. Narayanan, ``AI agents that matter,'' \emph{arXiv preprint arXiv:2407.01502}, 2024.

\bibitem{b_selfcorrect} J. Huang \emph{et al.}, ``Large language models cannot self-correct reasoning yet,'' in \emph{Proc. Int. Conf. Learning Representations (ICLR)}, 2024.

\bibitem{b_judge_survey} J. Gu \emph{et al.}, ``A survey on LLM-as-a-judge,'' \emph{arXiv preprint arXiv:2411.15594}, 2024.

\bibitem{b_length_bias} Y. Dubois, P. Liang, and T. Hashimoto, ``Length-controlled AlpacaEval: A simple way to debias automatic evaluators,'' in \emph{Proc. Conf. Language Modeling (COLM)}, 2024.

\bibitem{b_judge_inconsistent} R. Stureborg, D. Alikaniotis, and Y. Suhara, ``Large language models are inconsistent and biased evaluators,'' \emph{arXiv preprint arXiv:2405.01724}, 2024.

\bibitem{b_nondeterminism} S. Ouyang, J. M. Zhang, M. Harman, and M. Wang, ``An empirical study of the non-determinism of ChatGPT in code generation,'' \emph{ACM Transactions on Software Engineering and Methodology}, vol. 34, no. 2, pp. 1--28, 2025.

\bibitem{b_holm} S. Holm, ``A simple sequentially rejective multiple test procedure,'' \emph{Scandinavian Journal of Statistics}, vol. 6, no. 2, pp. 65--70, 1979.


% ===== Topic 1: Learning path / course sequence recommendation =====

\bibitem{b_shi2020kg} D. Shi \emph{et al.}, ``A learning path recommendation model based on a multidimensional knowledge graph framework for e-learning,'' \emph{Knowl.-Based Syst.}, vol. 195, Art. no. 105618, 2020, doi: 10.1016/j.knosys.2020.105618.
% topic: 1; claim: knowledge-graph-based learning path recommendation that links learning objects through several semantic relations to generate goal-specific paths; verified: https://api.crossref.org/works/10.1016/j.knosys.2020.105618

\bibitem{b_li2023gehrl} Q. Li \emph{et al.}, ``Graph enhanced hierarchical reinforcement learning for goal-oriented learning path recommendation,'' in \emph{Proc. 32nd ACM Int. Conf. Inf. Knowl. Manage. (CIKM)}, 2023, pp. 1318--1327, doi: 10.1145/3583780.3614897.
% topic: 1; claim: RL-based, goal-oriented learning path recommendation with a graph-based candidate selector that keeps paths restricted to goal-related items; verified: https://api.crossref.org/works/10.1145/3583780.3614897

\bibitem{b_amin2023rl} S. Amin \emph{et al.}, ``Smart e-learning framework for personalized adaptive learning and sequential path recommendations using reinforcement learning,'' \emph{IEEE Access}, vol. 11, pp. 89769--89790, 2023, doi: 10.1109/ACCESS.2023.3305584.
% topic: 1; claim: IEEE Access example of MDP/RL-based sequential learning path recommendation; verified: https://api.crossref.org/works/10.1109/access.2023.3305584

\bibitem{b_meng2021ldlp} L. Meng \emph{et al.}, ``LD--LP generation of personalized learning path based on learning diagnosis,'' \emph{IEEE Trans. Learn. Technol.}, vol. 14, no. 1, pp. 122--128, 2021, doi: 10.1109/TLT.2021.3058525.
% topic: 1; claim: IEEE TLT example of diagnosis-driven personalized learning path generation (prerequisite/knowledge-state-aware); verified: https://api.crossref.org/works/10.1109/tlt.2021.3058525

\bibitem{b_siren2022oer} A. Siren and V. Tzerpos, ``Automatic learning path creation using OER: A systematic literature mapping,'' \emph{IEEE Trans. Learn. Technol.}, vol. 15, no. 4, pp. 493--507, 2022, doi: 10.1109/TLT.2022.3193751.
% topic: 1; claim: systematic mapping of automatic learning path creation from open educational resources, i.e., the same select-and-order task over an open catalog; verified: https://api.crossref.org/works/10.1109/tlt.2022.3193751

\bibitem{b_wang2024hgcr} Y. Wang, D. Ma, J. Ma, and Q. Jin, ``HGCR: A heterogeneous graph-enhanced interactive course recommendation scheme for online learning,'' \emph{IEEE Trans. Learn. Technol.}, vol. 17, pp. 364--374, 2024, doi: 10.1109/TLT.2023.3314399.
% topic: 1; claim: graph-based (heterogeneous graph) course recommendation published in IEEE TLT; verified: https://api.crossref.org/works/10.1109/tlt.2023.3314399

\bibitem{b_liu2020prereq} L. Liu \emph{et al.}, ``Prerequisite relation learning for course concepts based on hyperbolic deep representation,'' \emph{IEEE Access}, vol. 8, pp. 49079--49089, 2020, doi: 10.1109/ACCESS.2020.2979555.
% topic: 1; claim: learning prerequisite relations between course concepts, which underpin prerequisite-aware ordering of learning paths; verified: https://api.crossref.org/works/10.1109/access.2020.2979555

\bibitem{b_bin2024study} Q. Bin, M. F. Zuhairi, and J. Morcos, ``A comprehensive study on personalized learning recommendation in e-learning system,'' \emph{IEEE Access}, vol. 12, pp. 100446--100482, 2024, doi: 10.1109/ACCESS.2024.3428419.
% topic: 1; claim: recent IEEE Access review of personalized learning recommendation techniques in e-learning; verified: https://api.crossref.org/works/10.1109/access.2024.3428419

\bibitem{b_mittal2026course} S. Mittal and I. Malhotra, ``Recent advancements in course recommendation systems in the era of large language models (2022--2025),'' \emph{IEEE Access}, vol. 14, pp. 53164--53188, 2026, doi: 10.1109/ACCESS.2026.3681211.
% topic: 1; claim: IEEE Access survey of LLM-era course recommendation (2022-2025), situating LLM-based curation within course recommender research; verified: https://api.crossref.org/works/10.1109/access.2026.3681211

% ===== Topic 2: RAG / LLMs for educational recommendation =====

\bibitem{b_li2025ragedu} Z. Li \emph{et al.}, ``Retrieval-augmented generation for educational application: A systematic survey,'' \emph{Comput. Educ.: Artif. Intell.}, vol. 8, Art. no. 100417, 2025, doi: 10.1016/j.caeai.2025.100417.
% topic: 2; claim: systematic survey of RAG in educational applications, supporting retrieval-grounded LLM pipelines as a standard design in education; verified: https://api.crossref.org/works/10.1016/j.caeai.2025.100417

% ===== Topic 3: Multi-agent LLMs in education; critiques; self-correction =====

\bibitem{b_yang2025agentsedu} J. Yang, M. Wang, X. Du, and R. Na, ``A comprehensive survey on large-language-model-based agents for education,'' \emph{IEEE Trans. Learn. Technol.}, vol. 18, pp. 898--913, 2025, doi: 10.1109/TLT.2025.3617909.
% topic: 3; claim: IEEE TLT survey of LLM-based (single- and multi-) agent systems for education; verified: https://api.crossref.org/works/10.1109/tlt.2025.3617909

\bibitem{b_nguyen2026mas} Q.-V. Nguyen, W.-Y. Hwang, M.-T. Sun, K. Sakai, and W.-S. Ku, ``LLM-based multi-agent system for personalized lesson plan generation in multi-domain adaptive learning,'' \emph{IEEE Trans. Learn. Technol.}, early access, 2026, doi: 10.1109/TLT.2026.3742134.
% topic: 3; claim: recent IEEE TLT example of a multi-agent LLM pipeline for personalized educational planning, the design family our study stress-tests; verified: https://api.crossref.org/works/10.1109/tlt.2026.3742134

\bibitem{b_cemri2025mas} M. Cemri \emph{et al.}, ``Why do multi-agent LLM systems fail?'' in \emph{Adv. Neural Inf. Process. Syst. (NeurIPS)}, vol. 38, 2025, pp. 135676--135729, doi: 10.52202/085713-4082.
% topic: 3; claim: taxonomy of failure modes of multi-agent LLM systems, showing gains over single-agent baselines are often small; verified: https://api.crossref.org/works/10.52202/085713-4082

\bibitem{b_kamoi2024selfcorrect} R. Kamoi, Y. Zhang, N. Zhang, J. Han, and R. Zhang, ``When can LLMs actually correct their own mistakes? A critical survey of self-correction of LLMs,'' \emph{Trans. Assoc. Comput. Linguist.}, vol. 12, pp. 1417--1440, 2024, doi: 10.1162/tacl\_a\_00713.
% topic: 3; claim: intrinsic self-correction with LLM-generated feedback rarely helps without reliable external feedback, motivating our critic-loop vs. programmatic-verifier comparison; verified: https://api.crossref.org/works/10.1162/tacl_a_00713

% ===== Topic 4: Reasoning models and System 1 / System 2 =====

\bibitem{b_kahneman2011} D. Kahneman, \emph{Thinking, Fast and Slow}. New York, NY, USA: Farrar, Straus and Giroux, 2011.
% topic: 4; claim: origin of the System 1 (fast, intuitive) vs. System 2 (slow, deliberative) framing used for reasoning-off vs. reasoning-on and for "System One" decision models; verified: https://openlibrary.org/isbn/9780374275631 (first publ. 2011, FSG)

\bibitem{b_ziabari2025spectrum} A. S. Ziabari \emph{et al.}, ``Reasoning on a spectrum: Aligning LLMs to System 1 and System 2 thinking,'' arXiv:2502.12470, 2025.
% topic: 4; claim: System-1-aligned LLMs can match or beat System-2 (step-by-step) variants on many tasks, so deliberate reasoning is not uniformly better; verified: https://arxiv.org/abs/2502.12470

\bibitem{b_zhang2026sys1sys2} D. Zhang \emph{et al.}, ``From System 1 to System 2: A survey of reasoning large language models,'' \emph{IEEE Trans. Pattern Anal. Mach. Intell.}, vol. 48, no. 3, pp. 3335--3354, 2026, doi: 10.1109/TPAMI.2025.3637037.
% topic: 4; claim: survey of reasoning ("System 2") LLMs such as o1 and DeepSeek-R1 (arXiv:2502.17419, where Z.-Z. Li is listed first); verified: https://api.crossref.org/works/10.1109/TPAMI.2025.3637037

\bibitem{b_chen2025overthinking} X. Chen \emph{et al.}, ``Do NOT think that much for 2+3=? On the overthinking of long reasoning models,'' in \emph{Proc. 42nd Int. Conf. Mach. Learn. (ICML)}, PMLR, vol. 267, 2025, pp. 9487--9499.
% topic: 4; claim: reasoning models spend excessive compute on easy problems with little accuracy gain (overthinking), relevant to reasoning_effort=high costs; verified: https://proceedings.mlr.press/v267/chen25bx.html

\bibitem{b_guo2025deepseekr1} D. Guo \emph{et al.}, ``DeepSeek-R1 incentivizes reasoning in LLMs through reinforcement learning,'' \emph{Nature}, vol. 645, no. 8081, pp. 633--638, 2025, doi: 10.1038/s41586-025-09422-z.
% topic: 4; claim: peer-reviewed description of an RL-trained long-chain-of-thought reasoning model, the System 2 model class; verified: https://api.crossref.org/works/10.1038/s41586-025-09422-z

\bibitem{b_openai2024o1} OpenAI, ``OpenAI o1 system card,'' arXiv:2412.16720, 2024.
% topic: 4; claim: system card of the first widely deployed reasoning model with test-time deliberation; verified: https://arxiv.org/abs/2412.16720

\bibitem{b_sprague2025cot} Z. Sprague \emph{et al.}, ``To CoT or not to CoT? Chain-of-thought helps mainly on math and symbolic reasoning,'' in \emph{Proc. Int. Conf. Learn. Represent. (ICLR)}, 2025. [Online]. Available: https://arxiv.org/abs/2409.12183
% topic: 4; claim: meta-analysis showing chain-of-thought gives large gains mainly on math/symbolic tasks and little elsewhere, consistent with small reasoning-on gains for curation; verified: https://arxiv.org/abs/2409.12183

% ===== Topic 5: Encoder / non-autoregressive decision models =====

\bibitem{b_warner2025modernbert} B. Warner \emph{et al.}, ``Smarter, better, faster, longer: A modern bidirectional encoder for fast, memory efficient, and long context finetuning and inference,'' in \emph{Proc. 63rd Annu. Meeting Assoc. Comput. Linguist. (ACL)}, 2025, pp. 2526--2547, doi: 10.18653/v1/2025.acl-long.127.
% topic: 5; claim: ModernBERT, the encoder backbone of the Laya English decision checkpoints; verified: https://api.crossref.org/works/10.18653/v1/2025.acl-long.127

\bibitem{b_marone2025mmbert} M. Marone \emph{et al.}, ``mmBERT: A modern multilingual encoder with annealed language learning,'' arXiv:2509.06888, 2025.
% topic: 5; claim: mmBERT, the multilingual encoder behind laya-multilingual (relevant for Portuguese input); verified: https://arxiv.org/abs/2509.06888

\bibitem{b_yin2019zeroshot} W. Yin, J. Hay, and D. Roth, ``Benchmarking zero-shot text classification: Datasets, evaluation and entailment approach,'' in \emph{Proc. Conf. Empirical Methods Natural Lang. Process. and 9th Int. Joint Conf. Natural Lang. Process. (EMNLP-IJCNLP)}, 2019, pp. 3912--3921, doi: 10.18653/v1/D19-1404.
% topic: 5; claim: NLI/entailment-based zero-shot classification with encoders, the precursor of non-generative decision models that score candidate labels; verified: https://api.crossref.org/works/10.18653/v1/d19-1404

\bibitem{b_zaratiana2024gliner} U. Zaratiana, N. Tomeh, P. Holat, and T. Charnois, ``GLiNER: Generalist model for named entity recognition using bidirectional transformer,'' in \emph{Proc. Conf. North Amer. Chapter Assoc. Comput. Linguist.: Hum. Lang. Technol. (NAACL-HLT)}, 2024, pp. 5364--5376, doi: 10.18653/v1/2024.naacl-long.300.
% topic: 5; claim: small bidirectional encoders can rival LLMs on open-label structured extraction while being much cheaper and faster; verified: https://api.crossref.org/works/10.18653/v1/2024.naacl-long.300

\bibitem{b_typesafe2026jev} D. Almeida, ``Introducing System One models \& Jev,'' TypeSafe AI, Sep. 15, 2026. [Online]. Available: https://typesafe.ai/blog/introducing-system-one-models-and-jev (accessed Oct. 2026).
% topic: 5; claim: official launch post of Jev, a hosted non-autoregressive "System One" model that returns typed decisions with calibrated probabilities, all sampled in parallel; verified: https://typesafe.ai/blog/introducing-system-one-models-and-jev

\bibitem{b_typesafe2026docs} TypeSafe AI, ``Introduction---Jev documentation,'' 2026. [Online]. Available: https://docs.typesafe.ai/introduction (accessed Oct. 2026).
% topic: 5; claim: API documentation describing Jev's state + typed-question interface (choice/score/boolean answers with probability distributions); verified: https://docs.typesafe.ai/introduction

\bibitem{b_palmer2026kev} J. Palmer, ``Kev: Jev-like decision models built on Qwen,'' GitHub repository, 2026. [Online]. Available: https://github.com/jaredpalmer/kev (accessed Oct. 2026).
% topic: 5; claim: open-source (Apache-2.0) Jev-like decision models on Qwen3.5/Qwen3.8 bases answering yes/no, multiple-choice and rating questions with calibrated probabilities; verified: https://github.com/jaredpalmer/kev

\bibitem{b_laya2026} NandhaKishorM (ConvAI Innovations), ``Laya: Non-autoregressive System 1 decision engine,'' GitHub repository, 2026. [Online]. Available: https://github.com/NandhaKishorM/laya (accessed Oct. 2026).
% topic: 5; claim: open-source (Apache-2.0) encoder-only decision engine (ModernBERT-large / mmBERT-base) answering typed choice/score/yes-no questions in one forward pass in 100+ languages; verified: https://github.com/NandhaKishorM/laya

% ===== Topic 6: LLM-as-a-judge reliability and agreement statistics =====

\bibitem{b_panickssery2024selfpref} A. Panickssery, S. R. Bowman, and S. Feng, ``LLM evaluators recognize and favor their own generations,'' in \emph{Adv. Neural Inf. Process. Syst. (NeurIPS)}, vol. 37, 2024, pp. 68772--68802, doi: 10.52202/079017-2197.
% topic: 6; claim: self-preference bias of LLM judges, motivating judges from two model families that differ from the generators; verified: https://api.crossref.org/works/10.52202/079017-2197

\bibitem{b_verga2024juries} P. Verga \emph{et al.}, ``Replacing judges with juries: Evaluating LLM generations with a panel of diverse models,'' arXiv:2404.18796, 2024.
% topic: 6; claim: a panel of diverse LLM judges reduces intra-model bias and correlates better with humans than one large judge; verified: https://arxiv.org/abs/2404.18796

\bibitem{b_liu2023geval} Y. Liu \emph{et al.}, ``G-Eval: NLG evaluation using GPT-4 with better human alignment,'' in \emph{Proc. Conf. Empirical Methods Natural Lang. Process. (EMNLP)}, 2023, pp. 2511--2522, doi: 10.18653/v1/2023.emnlp-main.153.
% topic: 6; claim: rubric-based LLM judging with chain-of-thought evaluation steps improves agreement with humans (basis for reason-first anchored rubric); verified: https://api.crossref.org/works/10.18653/v1/2023.emnlp-main.153

\bibitem{b_bavaresco2025judges} A. Bavaresco \emph{et al.}, ``LLMs instead of human judges? A large scale empirical study across 20 NLP evaluation tasks,'' in \emph{Proc. 63rd Annu. Meeting Assoc. Comput. Linguist. (ACL), Vol. 2: Short Papers}, 2025, pp. 238--255, doi: 10.18653/v1/2025.acl-short.20.
% topic: 6; claim: LLM-judge agreement with human annotators varies widely across tasks, so LLM judges should be validated before replacing humans; verified: https://api.crossref.org/works/10.18653/v1/2025.acl-short.20

\bibitem{b_hayes2007alpha} A. F. Hayes and K. Krippendorff, ``Answering the call for a standard reliability measure for coding data,'' \emph{Commun. Methods Meas.}, vol. 1, no. 1, pp. 77--89, 2007, doi: 10.1080/19312450709336664.
% topic: 6; claim: Krippendorff's alpha as the standard inter-rater reliability coefficient (ordinal metric, any number of raters); verified: https://api.crossref.org/works/10.1080/19312450709336664

\bibitem{b_cohen1968kappa} J. Cohen, ``Weighted kappa: Nominal scale agreement provision for scaled disagreement or partial credit,'' \emph{Psychol. Bull.}, vol. 70, no. 4, pp. 213--220, 1968, doi: 10.1037/h0026256.
% topic: 6; claim: weighted kappa for agreement between two raters on ordinal scales (inter-judge agreement); verified: https://api.crossref.org/works/10.1037/h0026256

% ===== Topic 7: Statistics =====

\bibitem{b_bates2015lme4} D. Bates, M. M{\"a}chler, B. Bolker, and S. Walker, ``Fitting linear mixed-effects models using lme4,'' \emph{J. Stat. Softw.}, vol. 67, no. 1, pp. 1--48, 2015, doi: 10.18637/jss.v067.i01.
% topic: 7; claim: linear mixed-effects models with crossed random effects (competency, judge, run); verified: https://api.crossref.org/works/10.18637/jss.v067.i01

\bibitem{b_efron1993bootstrap} B. Efron and R. J. Tibshirani, \emph{An Introduction to the Bootstrap}. New York, NY, USA: Chapman \& Hall, 1993, doi: 10.1007/978-1-4899-4541-9.
% topic: 7; claim: bootstrap confidence intervals for system-level score differences; verified: https://api.crossref.org/works/10.1007/978-1-4899-4541-9

\bibitem{b_wilcoxon1945} F. Wilcoxon, ``Individual comparisons by ranking methods,'' \emph{Biometrics Bull.}, vol. 1, no. 6, pp. 80--83, 1945, doi: 10.2307/3001968.
% topic: 7; claim: Wilcoxon signed-rank test for paired, non-normal per-competency comparisons; verified: https://api.crossref.org/works/10.2307/3001968

% ===== Topic 8: AI in the public sector / civil-servant training (incl. Brazil) =====

\bibitem{b_wirtz2019ai} B. W. Wirtz, J. C. Weyerer, and C. Geyer, ``Artificial intelligence and the public sector---Applications and challenges,'' \emph{Int. J. Public Admin.}, vol. 42, no. 7, pp. 596--615, 2019, doi: 10.1080/01900692.2018.1498103.
% topic: 8; claim: overview of AI applications and challenges in the public sector; verified: https://api.crossref.org/works/10.1080/01900692.2018.1498103

\bibitem{b_zuiderwijk2021gov} A. Zuiderwijk, Y.-C. Chen, and F. Salem, ``Implications of the use of artificial intelligence in public governance: A systematic literature review and a research agenda,'' \emph{Gov. Inf. Q.}, vol. 38, no. 3, Art. no. 101577, 2021, doi: 10.1016/j.giq.2021.101577.
% topic: 8; claim: systematic review of AI in public governance calling for empirical, context-specific evaluations; verified: https://api.crossref.org/works/10.1016/j.giq.2021.101577

\bibitem{b_ribeiro2025rap} M. M. Ribeiro and C. I. Segatto, ``Intelig{\^e}ncia artificial nas organiza{\c{c}}{\~o}es p{\'u}blicas brasileiras: Heterogeneidades e capacidades em tecnologia da informa{\c{c}}{\~a}o,'' \emph{Rev. Admin. P{\'u}blica}, vol. 59, no. 1, Art. no. e2024-0066, 2025, doi: 10.1590/0034-761220240066.
% topic: 8; claim: Brazilian evidence that AI adoption in public organizations is heterogeneous and tied to IT capacity (TIC Governo Eletronico 2021 microdata); verified: https://api.crossref.org/works/10.1590/0034-761220240066

\bibitem{b_santos2025plafor} T. D. dos Santos, L. L. S. Silva, and M. A. de Azevedo, ``Plano de Forma{\c{c}}{\~a}o Continuada dos Servidores da Rede Federal de Educa{\c{c}}{\~a}o Profissional, Cient{\'i}fica e Tecnol{\'o}gica e seus caminhos at{\'e} a plataformiza{\c{c}}{\~a}o,'' \emph{Rev. Educ. Pol{\'i}t. Debate}, vol. 14, no. 2, pp. 1--18, 2025, doi: 10.14393/REPOD-v14n2a2025-74251.
% topic: 8; claim: history of the Plafor continuing-training plan for Brazilian federal network civil servants and its move to a digital platform (PlaforEdu context); verified: https://seer.ufu.br/index.php/revistaeducaopoliticas/article/view/74251

\bibitem{b_ribeiro2022digcompedu} L. O. M. Ribeiro and L. C. G. Silveira, ``DigCompEdu e a plataforma digital de forma{\c{c}}{\~a}o de servidores: O uso do framework te{\'o}rico europeu nas compet{\^e}ncias digitais para profissionais da educa{\c{c}}{\~a}o da Rede Federal EPCT,'' \emph{EmRede -- Rev. Educ. Dist.}, vol. 9, no. 2, 2022, doi: 10.53628/emrede.v9i2.893.
% topic: 8; claim: competency framework (DigCompEdu) used to organize the federal civil-servant training platform, supporting the competency-based framing of learning paths; verified: https://api.crossref.org/works/10.53628/emrede.v9i2.893

\end{thebibliography}
